\section{Next-token Prediction: cómo un solo objetivo abarca una enorme cantidad de tareas}

La sección anterior presentó el método; esta sección aborda la primera intuición central de Jason. En apariencia, el entrenamiento de un modelo de lenguaje solo pide predecir el siguiente token, pero la probabilidad correcta de cada token puede depender de la sintaxis, el significado de las palabras, los hechos, las relaciones espaciales, la traducción, la aritmética e incluso el género del documento. Lo importante no es ponerle al preentrenamiento la etiqueta «multi-task», sino entender cómo el contexto forma de manera temporal los límites de cada tarea.

\subsection{De la distribución de probabilidad a la pérdida de entrenamiento}

Esta subsección completa primero el trasfondo matemático mínimo. Dados los tokens previos $x_{<t}=(x_1,\ldots,x_{t-1})$, el modelo produce la probabilidad condicional de cada token candidato del vocabulario; durante el entrenamiento solo se observa el token real $x_t$, y se aumenta su probabilidad. Al leer la slide, conviene distinguir primero entre «el modelo produce la distribución completa» y «los datos solo aportan un resultado observado».

\lecturefigure{jason-slide-004.jpg}{El modelo de lenguaje asigna a cada palabra candidata del vocabulario una probabilidad para la siguiente posición}{Jason Wei official deck, p.~4}

La log-verosimilitud negativa de un modelo de lenguaje autorregresivo puede escribirse como
\[
\mathcal{L}_{\mathrm{NLL}}(\theta)=-\frac{1}{T}\sum_{t=1}^{T}\log p_{\theta}(x_t\mid x_{<t}).
\]
Aquí $\theta$ representa los parámetros del modelo, $T$ es la longitud de la secuencia y $p_{\theta}(x_t\mid x_{<t})$ es la probabilidad que el modelo asigna al siguiente token real. Cuanto menor es la probabilidad, mayor es la penalización $-\log p$; el optimizador actualiza $\theta$ mediante el gradiente para que la continuation real de la distribución de entrenamiento sea más probable.

\begin{knowledgebox}{Qué es la perplexity}
La perplexity (perplejidad, PPL) es la exponencial de la entropía cruzada promedio:
\[
\operatorname{PPL}=\exp\!\left(\mathcal{L}_{\mathrm{NLL}}\right).
\]
Puede entenderse de manera aproximada como el «número efectivo de candidatos» que el modelo enfrenta en cada paso. Por ejemplo, una PPL de 10 no significa que el modelo elija realmente de manera uniforme entre diez palabras; es solo un resumen exponencial de la pérdida logarítmica, y tampoco equivale directamente a la exactitud factual, la tasa de aciertos en razonamiento ni las preferencias humanas.
\end{knowledgebox}

\subsection{El massive multi-task learning no es una metáfora}

Con la pérdida ya definida, el siguiente paso es preguntar qué determina la probabilidad de un token. Si el corpus es lo bastante grande, el modelo encontrará una y otra vez, en contextos distintos, sintaxis, semántica, hechos y formatos de tarea. Cada regularidad predecible aporta gradiente a la misma NLL; por ello, el objetivo unificado acumula en el espacio de parámetros muchas señales de aprendizaje locales.

\lecturefigure{jason-slide-005.jpg}{Intuición uno: la next-word prediction a gran escala es massive multi-task learning}{Jason Wei official deck, p.~5}

\begin{importantbox}{Definición de latent task}
Una latent task (tarea latente) es un subproblema de predicción que el contexto define de manera temporal, sin un task ID explícito. Por ejemplo, «The capital of Azerbaijan is» activa la recuperación de un hecho, «The movie was» puede activar una compleción de sentimiento y «The Spanish word for pretty is» activa una traducción. Comparten un mismo vocabulario y una misma pérdida, pero requieren cálculos intermedios distintos.
\end{importantbox}

\lecturefigure{jason-slide-006.jpg}{Las oraciones de preentrenamiento enseñan a la vez sintaxis, semántica, hechos, sentimiento, traducción, relaciones espaciales y aritmética}{Jason Wei official deck, p.~6}

\begin{knowledgebox}{Lectura de la figura: comparar primero «de dónde viene la supervisión»}
Ninguna fila etiqueta por separado el nombre de la tarea; la supervisión proviene de la continuation real. El ejemplo de sintaxis exige descartar palabras que no encajan sintácticamente; el de conocimiento del mundo exige colocar Baku por delante de London; el espacial exige seguir la posición de los personajes; el aritmético exige aprovechar las relaciones del enunciado. Lo que la figura respalda no es que «el modelo aprenda con seguridad cada capacidad», sino que «el corpus efectivamente aporta gradientes de probabilidad hacia estas capacidades».
\end{knowledgebox}

\teachervoice{Entre 00:04:59 y 00:07:03, Jason presenta varios ejemplos seguidos y al final enfatiza que puede haber millones de tareas así. El tono de la clase es importante: no afirma que los datos de preentrenamiento estén equilibrados de forma natural, sino que, mientras existan las oraciones pertinentes, el next-token objective puede extraer supervisión de ellas.}

\subsection{Las tareas pueden ser tan arbitrarias que resulta difícil nombrarlas}

Las tareas de la página anterior son todas muy «limpias» y fáciles de nombrar para una persona; esta subsección examina a propósito una continuation difícil de clasificar. Si una oración de Wikipedia pide predecir, en orden, un nombre de persona, una coma, un artículo y una palabra de profesión, cada posición recurre a una regularidad distinta. Las tareas no son un benchmark enumerado de antemano, sino preguntas momentáneas que plantea cada distribución condicional de la secuencia.

\lecturefigure{jason-slide-007.jpg}{Las tareas latentes de una misma oración pueden saltar de un hecho a la puntuación, y luego a regularidades de continuación difíciles de nombrar}{Jason Wei official deck, p.~7}

Si el token $t$-ésimo del corpus implica una tarea $z_t$, el objetivo de entrenamiento puede escribirse conceptualmente como
\[
\mathcal{L}=\mathbb{E}_{(x,z)\sim\mathcal{D}}\bigl[\ell(x_t,x_{<t};z_t)\bigr].
\]
Aquí $\mathcal{D}$ es la distribución de entrenamiento; $z_t$ no es una etiqueta que los datos proporcionen de forma explícita, sino una variable latente que introduce el analista para facilitar la comprensión, y $\ell$ es la pérdida de predicción en esa posición. Esta notación recuerda que un mismo token puede depender de varias regularidades a la vez; por eso, la descomposición en tareas es solo una herramienta explicativa y no garantiza que exista un límite único y correcto.

\begin{warningbox}{No hay que imaginar el pretraining corpus como un almacén de tareas ya depurado automáticamente}
El mismo mecanismo también aprende errores ortográficos, sesgos, correlaciones espurias, copias de plantillas y fuentes de baja calidad. En la sesión de Q\&A, Jason dice explícitamente que el entrenamiento no distingue de forma automática entre datos buenos y malos; los investigadores todavía deben revisar el origen de los datos, filtrar el contenido claramente poco confiable y comprobar, mediante evaluaciones independientes, si lo aprendido son regularidades transferibles o atajos de los datos.
\end{warningbox}

\begin{lstlisting}[caption={Pseudocódigo para descomponer un texto en latent-task probes}]
for position in document:
    context = tokens[:position]
    target = tokens[position]
    ask("What knowledge or rule raises target probability?")
    tag(grammar, fact, format, reasoning, noise, other)
    store(context, target, tags)
\end{lstlisting}

\subsection{Resumen de la sección}

La fuerza de la next-token prediction no se debe a que el nombre del objetivo sea complejo, sino a que el corpus reúne una gran cantidad de problemas de predicción condicional en un mismo objetivo. En cada paso, el modelo solo optimiza la probabilidad del token real, pero acumula en sus parámetros compartidos representaciones y patrones de cálculo que atraviesan distintas tareas. Esta perspectiva prepara la sección siguiente: si la loss global es una mezcla de las pérdidas de muchas tareas latentes, la curva global y las curvas de cada capacidad no tienen por qué tener la misma forma.
